\chapter{Guia Practica: Como Utilizar el Dataset}
\label{cap:usage_guide}

\section{Acceso al Dataset}

El dataset FatigueSet esta disponible en:
\url{https://www.esense.io/datasets/fatigueset}

Descargue el archivo completo comprimido (~500 MB aproximadamente).

\section{Estructura de Directorios Despues de Descargar}

\begin{verbatim}
fatigueset/
├── README.md              (Especificaciones - LEA PRIMERO)
├── metadata.csv           (Contrabalanceo de participantes)
├── 01/ ... 12/            (Directorios de participantes)
│   ├── 01/ 02/ 03/        (Sesiones 1, 2, 3)
│   │   ├── exp_crt.csv
│   │   ├── exp_nback.csv
│   │   ├── exp_task_switch.csv
│   │   ├── exp_fatigue.csv
│   │   ├── exp_markers.csv
│   │   ├── ear_*.csv (6 archivos)
│   │   ├── forehead_*.csv (13 archivos)
│   │   ├── chest_*.csv (9 archivos)
│   │   └── wrist_*.csv (6 archivos)
└── fatigueset_aggregated_features.csv  (ARCHIVO PRINCIPAL PARA ANALISIS)
\end{verbatim}

\section{Archivo Principal para Analisis: fatigueset\_aggregated\_features.csv}

Para 95\% de los analisis, use este archivo en lugar de los 34 CSV crudos por sesion.

\subsection{Como Cargar en Python}

\begin{lstlisting}[language=Python]
import pandas as pd
import numpy as np

# Cargar datos
df = pd.read_csv('fatigueset_aggregated_features.csv')

# Inspeccionar
print(f"Forma: {df.shape}")  # 108 filas, 19 columnas
print(df.head())
print(df.describe())

# Limpiar valores problematicos
df_clean = df.replace([np.inf, -np.inf], np.nan)
df_clean = df_clean.dropna()
print(f"Registros validos: {len(df_clean)}")
\end{lstlisting}

\section{Tareas Comunes}

\subsection{1. Comparar Variables entre Fases}

\begin{lstlisting}[language=Python]
# Agrupar por fase
for fase in ['M1_baseline', 'M2_post_ejercicio', 'M3_post_fatiga_mental']:
    fase_data = df_clean[df_clean['fase'] == fase]
    print(f"\n{fase}:")
    print(f"  HR: {fase_data['hr_media'].mean():.1f} bpm")
    print(f"  HRV: {fase_data['hrv_media'].mean():.1f} ms")
    print(f"  EDA: {fase_data['eda_media'].mean():.2f} uS")
\end{lstlisting}

\subsection{2. Analizar Efecto de Intensidad}

\begin{lstlisting}[language=Python]
# Comparar low vs medium vs high
for intensidad in ['low', 'medium', 'high']:
    int_data = df_clean[df_clean['intensidad'] == intensidad]
    print(f"\n{intensidad.upper()}:")
    print(f"  HR promedio: {int_data['hr_media'].mean():.1f}")
    print(f"  HRV promedio: {int_data['hrv_media'].mean():.1f}")
    print(f"  Fatiga Fisica reportada: {int_data['fatiga_fisica'].mean():.1f}/100")
\end{lstlisting}

\subsection{3. Visualizaciones Basicas}

\begin{lstlisting}[language=Python]
import matplotlib.pyplot as plt
import seaborn as sns

# Box plot de HR por fase
sns.boxplot(data=df_clean, x='fase', y='hr_media')
plt.title('Frecuencia Cardiaca por Fase')
plt.show()

# Scatter plot: HRV vs Autoreporte Fatiga
plt.scatter(df_clean['hrv_media'], df_clean['fatiga_fisica'], alpha=0.6)
plt.xlabel('HRV (ms)')
plt.ylabel('Fatiga Fisica Reportada')
plt.title('Relacion HRV vs Autoreporte Fatiga')
plt.show()
\end{lstlisting}

\section{Notas Importantes}

\subsection{Valores Problematicos}

El dataset puede contener:
\begin{itemize}
    \item \textbf{-inf, inf, NaN}: Generalmente participante 12, algunas fases.
          Solucion: usar `dropna()` o `replace([np.inf, -np.inf], np.nan)`
    \item \textbf{Valores atipicos}: Algunos participantes pueden tener respuestas extremas.
          Normalice por participante si es necesario.
\end{itemize}

\subsection{Unidades y Escalas}

\begin{table}[H]
\centering
\caption{Unidades para cada variable}
\begin{tabular}{|l|l|l|}
\hline
\textbf{Variable} & \textbf{Unidad} & \textbf{Rango Tipico} \\
\hline
hr\_media & bpm & 60-120 \\
br\_media & brpm & 12-30 \\
hrv\_media & ms & 20-100 \\
eda\_media & µS & 0.1-3 \\
eeg\_*\_media & Bels & 0-1 \\
fatiga\_fisica, fatiga\_mental & escala 0-100 & 0-100 \\
\hline
\end{tabular}
\end{table}

\subsection{Contexto Temporal}

Este es un dataset de \textbf{laboratorio controlado}. No es datos "salvajes" de personas
en vida real. Protocolos:
\begin{itemize}
    \item Actividad fisica controlada (no sabemos exactamente que nivel de ejercicio)
    \item Tareas cognitivas estandarizadas
    \item Ambiente controlado (sin distracciones externas)
\end{itemize}

Generalizacion a escenarios reales requiere caucion.

\section{Proximos Pasos Sugeridos}

Para investigadores interesados en este dataset:

\begin{enumerate}
    \item \textbf{Validar resultados}: Reproduzca los analisis mostrados en Cap. 6
    \item \textbf{Analisis avanzado}: Modelos ML para prediccion de fatiga
    \item \textbf{Individualizacion}: Desarrolle modelos personalizados por sujeto
    \item \textbf{Series temporales}: Analice dinamica temporal dentro de sesiones
    \item \textbf{Comparaciones}: Compare contra otros datasets de fatiga si disponibles
\end{enumerate}

\section{Contacto y Referencia}

\textbf{Para problemas con el dataset}:
Consulte el archivo README.md en la raiz del directorio.

\textbf{Para citar este dataset}:
Utilice la referencia bibliografica de Kalanadhabhatta et al. (2021) en References.

\textbf{Para reportar errores/bug}s:
Contacte a los autores a traves del sitio web de eSense.
